Six flights of one airframe, each isolating one knob: which atmosphere
supplies the air, which integrator does the arithmetic, how big the timestep
is, and which way the launch rail points. The airframe is Tutorial~2's 26\,mm
bird on a D12; if you skipped ahead, the script rebuilds it.

\begin{note}[Setup]
An empty directory with \code{estes\_small.qmd} copied in;
\srcf{scripts/tutorial4.cli} replays the whole study. Each \repl{launch} is
followed by a \repl{save} so its trajectory survives the next flight.
\end{note}

\subsection*{The knobs}

\begin{console}
\muted{# QtRocket CLI - type 'help' for commands, 'quit' to exit.}
\uin{loaddb estes_small.qmd}
\ok loaddb: 13 new motors from estes_small.qmd (13 total)
\uin{newdesign NoseCone name=nose baseRadius=0.013 length=0.09 density=900}
\ok newdesign: root NoseCone id=2
\uin{addpart nose BodyTube name=body innerRadius=0.0121 outerRadius=0.013 length=0.30 density=1000}
\ok addpart: BodyTube id=3 under 2
\uin{addpart body FinSet name=fins finCount=3 rootChord=0.05 tipChord=0.025 span=0.03 sweep=0.02 thickness=0.003 bodyRadius=0.013 density=900 seat=OnSurface parentStation=0 childStation=0}
\ok addpart: FinSet id=4 under 3
\uin{checkdesign}
\ok checkdesign: 3 parts, contiguous, no overlaps
\uin{setmotor D12}
\ok setmotor: D12
\uin{setdrag 0.75}
\ok setdrag: 0.75
\uin{setarea 0.000531}
\ok setarea: 0.000531 m^2
\uin{listatmospheres}
\ok listatmospheres: 3 available
Constant Atmosphere
US Standard 1976
Vacuum
\uin{listgravity}
\ok listgravity: 2 available
Constant Gravity
Spherical Gravity
\uin{listintegrators}
\ok listintegrators: 2 available
Runge-Kutta 4th Order
Runge-Kutta-Fehlberg
\end{console}


Three \repl{list...} commands enumerate what is pluggable: three atmosphere
models, two gravity models, two integrators. \repl{setatmosphere},
\repl{setgravity}, and \repl{setintegrator} accept exactly these names ---
misspell one and the CLI rejects it and keeps the current model.

\subsection*{Baseline: constant air, RK4, $\Delta t = 0.01$\,s}

\begin{console}
\uin{settimestep 0.01}
\ok settimestep: 0.01 s
\uin{status}
\ok status:
  motor      = D12
  dry_mass   = 1 kg
  drag_coeff = 0.75
  ref_area   = 0.000531 m^2
  velocity   = 0 m/s
  angle      = 0 deg (from vertical)
  atmosphere = Constant Atmosphere
  gravity    = Constant Gravity
  integrator = Runge-Kutta 4th Order
  database   = 13 motors
\uin{launch}
\ok launch: 2037 steps, t_final=20.360s
  apogee    = 447.674 m @ t=8.140s
  max_speed = 145.186 m/s @ t=1.590s
  downrange = 0.000 m
  landing   = t=20.360s, pos=(0.000, 0.000, -0.350)
CSV /tmp/flight1/qtrocket_run.csv
\uin{save baseline_constant.csv}
\ok save: /tmp/flight1/baseline_constant.csv
\end{console}


$447.7\m$. Drag is computed from the \emph{active atmosphere's} density at the
rocket's current altitude --- that is the whole meaning of the atmosphere
setting. \code{Constant Atmosphere} uses sea-level density everywhere.

\subsection*{Real air: US Standard 1976}

\begin{console}
\uin{setatmosphere US Standard 1976}
\ok setatmosphere: US Standard 1976
\uin{launch}
\ok launch: 2048 steps, t_final=20.470s
  apogee    = 452.936 m @ t=8.210s
  max_speed = 145.514 m/s @ t=1.600s
  downrange = 0.000 m
  landing   = t=20.470s, pos=(0.000, 0.000, -0.226)
CSV /tmp/flight1/qtrocket_run.csv
\uin{save us_standard.csv}
\ok save: /tmp/flight1/us_standard.csv
\end{console}


$452.9\m$ --- about one percent higher, because air thins with altitude and a
450-meter flight barely samples the difference. Model fidelity buys little on
low flights and progressively more as you climb.

\subsection*{No air at all: the vacuum ceiling}

\begin{console}
\uin{setatmosphere Vacuum}
\ok setatmosphere: Vacuum
\uin{launch}
\ok launch: 4338 steps, t_final=43.370s
  apogee    = 2223.647 m @ t=22.070s
  max_speed = 208.924 m/s @ t=43.370s
  downrange = 0.000 m
  landing   = t=43.370s, pos=(0.000, 0.000, -1.841)
CSV /tmp/flight1/qtrocket_run.csv
\uin{save vacuum_rk4.csv}
\ok save: /tmp/flight1/vacuum_rk4.csv
\end{console}


$2223.6\m$ --- five times the apogee, from the same motor. Vacuum reduces the
model to thrust plus gravity, which makes it two useful things: a
deterministic upper bound (``what is drag costing me?'' --- here, 80\% of the
theoretical altitude) and a diagnostic (if a flight surprises you, fly it in
vacuum to separate propulsion questions from aerodynamics questions).

\subsection*{Same flight, adaptive integrator}

\begin{console}
\uin{setintegrator Runge-Kutta-Fehlberg}
\ok setintegrator: Runge-Kutta-Fehlberg
\uin{launch}
\ok launch: 528 steps, t_final=43.290s
  apogee    = 2223.530 m @ t=21.990s
  max_speed = 209.023 m/s @ t=43.290s
  downrange = 0.000 m
  landing   = t=43.290s, pos=(0.000, 0.000, -4.080)
CSV /tmp/flight1/qtrocket_run.csv
\uin{save vacuum_rkf45.csv}
\ok save: /tmp/flight1/vacuum_rkf45.csv
\end{console}


The Runge--Kutta--Fehlberg integrator adapts its own step size against an
error estimate. Its vacuum apogee agrees with fixed-step RK4 to within
$0.12\m$ in $2223$ --- five parts in a hundred thousand --- while choosing its
own steps. Agreement between two independent integration schemes is the
cheapest confidence check numerical work offers.

\subsection*{What a coarse timestep costs}

\begin{console}
\uin{setintegrator Runge-Kutta 4th Order}
\ok setintegrator: Runge-Kutta 4th Order
\uin{settimestep 0.1}
\ok settimestep: 0.1 s
\uin{launch}
\ok launch: 430 steps, t_final=42.900s
  apogee    = 2183.363 m @ t=21.800s
  max_speed = 207.113 m/s @ t=42.900s
  downrange = 0.000 m
  landing   = t=42.900s, pos=(0.000, 0.000, -3.702)
CSV /tmp/flight1/qtrocket_run.csv
\uin{save coarse_dt.csv}
\ok save: /tmp/flight1/coarse_dt.csv
\uin{settimestep 0.01}
\ok settimestep: 0.01 s
\end{console}


Back on RK4 with $\Delta t = 0.1$\,s --- ten times coarser --- the same vacuum
flight reads $2183.4\m$: forty meters (1.8\%) short, lost to integration error
across the sharp thrust transient.

\subsection*{Tilt the rail}

\begin{console}
\uin{setatmosphere US Standard 1976}
\ok setatmosphere: US Standard 1976
\uin{setangle 10}
\ok setangle: 10 deg
\uin{setvelocity 15}
\ok setvelocity: 15 m/s
\uin{launch}
\ok launch: 2106 steps, t_final=21.050s
  apogee    = 479.803 m @ t=8.280s
  max_speed = 151.030 m/s @ t=1.590s
  downrange = 12.127 m
  landing   = t=21.050s, pos=(12.127, 0.000, -0.118)
CSV /tmp/flight1/qtrocket_run.csv
\uin{save angled.csv}
\ok save: /tmp/flight1/angled.csv
\uin{quit}
\ok bye
\end{console}


Two changes together: a $15$\,m/s rail-exit velocity (\repl{setvelocity}) and
a $10^\circ$ tilt from vertical (\repl{setangle}). For the first time
\code{downrange} is nonzero --- $12.1\m$ of drift --- and apogee ticks up
relative to the upright US-Standard flight because the rail speed adds kinetic
energy. \Cref{fig:tut4-plots} shows both studies: the three atmosphere curves
overlaid, and the tilted flight's ground track.

\begin{figure}[htbp]\centering
\begin{subfigure}[b]{0.55\linewidth}\centering
\begin{tikzpicture}
\begin{axis}[qtplot, width=\linewidth, height=0.75\linewidth,
  xlabel={time (s)}, ylabel={altitude (m)}, xmin=0, ymin=0,
  legend pos=north east]
\addplot table {data/tut4-constant.dat}; \addlegendentry{Constant}
\addplot table {data/tut4-usstd.dat};    \addlegendentry{US Std 1976}
\addplot table {data/tut4-vacuum.dat};   \addlegendentry{Vacuum}
\end{axis}
\end{tikzpicture}
\caption{Same rocket, three atmospheres.}
\end{subfigure}\hfill
\begin{subfigure}[b]{0.42\linewidth}\centering
\begin{tikzpicture}
\begin{axis}[qtplot, width=\linewidth, height=0.98\linewidth,
  xlabel={downrange $x$ (m)}, ylabel={altitude (m)}, ymin=0]
\addplot table {data/tut4-angled.dat};
\end{axis}
\end{tikzpicture}
\caption{Ground track, $10^\circ$ tilt.}
\end{subfigure}
\caption{Tutorial 4's studies, plotted from the \repl{save}d CSVs: the
atmosphere overlay (left) and the tilted launch's $x$--$z$ trajectory
(right).}
\label{fig:tut4-plots}
\end{figure}

\begin{tip}[What you learned]
Atmosphere selection is $\rho(z)$ selection; vacuum is a bound and a
diagnostic; two integrators agreeing is cheap confidence; coarse timesteps
have a measurable price to pay for speed; and the launch rail's speed and angle are the last
two \repl{status} rows you had not touched yet.
\end{tip}
